\chapter{Introduction}

\section{Context and Motivation}

The landscape of technical education has undergone considerable transformation
in recent years, driven in large part by the proliferation of artificial
intelligence tools capable of generating content on demand. Within computing
and software engineering education, the ability to assess understanding and
reinforce learning through practice problems, quizzes, and coding challenges
has always been central to how students develop competency. However, the
process of creating these materials has traditionally been a time consuming
and manual endeavour, placing a significant burden on educators and leaving
students who wish to practise independently with a limited and often static
pool of exercises.

This project addresses that gap directly. It presents a full stack web
application, informally referred to during development as Code Quiz, that harnesses the generative capabilities of large language
models \cite{brown2020gpt3} to produce both multiple choice questions and
coding challenges on demand, based on a topic or prompt provided by the user.
Rather than returning a raw block of AI generated text, the platform enriches
every question it produces with two additional layers of machine learning
intelligence: a predicted set of algorithmic topic tags, and a predicted
difficulty level. These enrichments are produced by two custom trained
classifiers that operate on the text of each generated question, running
silently in the background without adding meaningful latency to the user
experience.

The motivation for building such a system sits at the intersection of three
well established problems. First, static question banks grow stale and do
not adapt to the breadth or specificity of what a learner wishes to study.
A student who wants to practise tree traversal problems on a Tuesday evening
should not be limited to the same twenty or thirty problems that a textbook
happens to cover, and an instructor who wishes to set coursework on a newly
introduced topic should not have to author every problem by hand. Second,
even when generic AI tools such as ChatGPT are used to generate questions,
the output is typically unstructured and requires manual curation before it
can be deployed as a teaching resource. The questions come back as free text
without structured options, without test cases, without tags, and without any
indication of how challenging they are likely to be. Third, learners
attempting self directed study with coding problems often have no objective
signal about how difficult a problem is or what underlying concept it is
testing, making it difficult to build a structured learning path from
simpler problems to harder ones.

By combining generative AI with trained machine learning models, this
platform attempts to address all three problems at once: questions are
produced dynamically to any specification, they are automatically tagged
with the algorithmic topics they assess, and they are labelled with a
predicted difficulty so that learners can make informed choices about which
problems to attempt. Crucially, the topic tagger and the difficulty predictor
are not wrappers around a second call to a language model. They are
lightweight scikit learn pipelines \cite{pedregosa2011} trained on real data:
a LeetCode problem corpus for the tag classifier \cite{leetcode2025} and the
IBM Project CodeNet dataset for the difficulty classifier \cite{puri2021codenet}.
This decision was deliberate. Running a second LLM call per question would
be slow, expensive, and would provide no reproducibility from one run to the
next, whereas a serialised classifier loaded once at startup produces
deterministic annotations in a few milliseconds per question.

\section{Project Scope and Level}

This is a level 8 project developed as a final year applied project in
computing. It encompasses a range of technically demanding areas, including
full stack web development, integration with external AI APIs, local model
inference, custom machine learning model training and deployment,
containerised infrastructure, and a social feature set that includes a
user accounts system with experience points and a friends network.

The scope is deliberately broad. The intent is not to build a research
prototype that solves a single narrowly defined problem, but rather to
deliver a functional, end to end product that a real user could interact with.
The application is fully containerised using Docker     \cite{merkel2014docker},
meaning it can be deployed in a reproducible way with a single command. The
frontend is a single page React application \cite{react2013} served through
Nginx \cite{nginx2025}. The backend is a Python FastAPI server
\cite{fastapi2025} that coordinates AI generation, machine learning inference,
and code execution. A PostgreSQL database \cite{postgres2025} persists user
data. The two trained classifiers are serialised as model artifacts and
loaded at server startup so that inference imposes no disk read cost on
request handling.

The project does not attempt to be a drop in replacement for professional
learning management systems, nor does it compete with platforms such as
LeetCode at scale. Its contribution is in demonstrating that a tightly
integrated pipeline connecting generative AI with custom trained classifiers
can enrich automatically produced educational content in a meaningful and
technically sound way, and that this can be achieved within the constraints
of a solo final year project using open source tooling throughout.

\section{Objectives}

The following objectives guided the development of this project and will
each be addressed explicitly in the system evaluation chapter:

\begin{enumerate}
    \item Design and implement a full stack web application that allows
    users to generate multiple choice and coding challenge quizzes through
    a natural language prompt, with support for multiple AI backends,
    including a cloud hosted model (OpenAI) and a locally hosted model
    (Ollama running Llama 3.1).
    \item Train a multi label topic tag classifier capable of predicting
    relevant algorithmic topic tags for a given coding question, using a
    real world dataset of labelled programming problems.
    \item Train a difficulty classifier capable of predicting whether a
    coding problem is easy, medium, or hard, in the absence of explicit
    human difficulty labels, by engineering a meaningful proxy label from
    available submission data.
    \item Integrate both trained classifiers into the backend inference
    pipeline so that every generated question is automatically enriched
    with predicted tags and a predicted difficulty level before being
    returned to the user.
    \item Implement a code execution engine within the backend capable of
    running user submitted code against test cases across multiple
    programming languages, and returning pass or fail results per test.
    \item Deliver a user accounts system with authentication, profile
    management, an experience point and levelling mechanic, and a social
    friends network.
\end{enumerate}

The primary metric for evaluating the machine learning components is
classification performance on held out test sets, using precision, recall,
and the F1 score. For the topic classifier specifically, F1 micro and F1
macro will both be reported, as they each reveal different aspects of
performance when class imbalance is present: F1 micro is dominated by the
most populous classes, whereas F1 macro treats every class equally regardless
of support. For the difficulty classifier, F1 macro is the primary metric
given the three class setup with roughly balanced support. For the
application itself, the evaluation focuses on whether the pipeline functions
correctly, whether generated questions receive sensible tags and difficulty
labels, and whether the code execution engine correctly judges test case
outcomes across every supported language.

\section{Chapter Overview}

The remainder of this dissertation is structured as follows. Chapter
2 describes the development approach, the planning and requirements
work done upfront, the GitHub based workflow, the continuous
integration pipeline, and the validation strategies for both software
and machine learning components. Chapter 3 surveys the technologies,
libraries, frameworks, and models that underpin the system, covering
large language models, the machine learning techniques used in both
classifiers, the web technologies used for the frontend and backend,
containerisation, and the training datasets. Chapter 4 presents the
architecture of the application, including component diagrams, data
flow descriptions, the database schema, the full training pipelines,
and the design of the inference bridge between the backend and the
trained model artefacts. Chapter 5 assesses the project against the
objectives, presenting classification metrics, confusion matrices, per
class breakdowns, and an honest examination of limitations. Chapter 6
concludes with a summary of outcomes, unexpected insights that
emerged, and concrete opportunities for future investigation.

\section{Project Repository}

The full source code, including the frontend, backend, training
scripts, Docker configuration, and serialised model artefacts, is
available at the following repository:

\begin{center}
\url{https://github.com/fionnmcgoldrick123/programming-quiz-app}
\end{center}

A brief screencast demonstrating the application is also available via
the repository README.

The repository contains a \texttt{frontend} directory (React and
TypeScript source), a \texttt{backend} directory (FastAPI application
with route handlers, services, auth, and the code execution engine),
an \texttt{ml\_models} directory (training scripts and serialised
\texttt{.pkl} artefacts for both classifiers), a
\texttt{docker-compose.yml} at the root that brings all three services
up in a single step, and a \texttt{docs} directory containing the
LaTeX source for this dissertation.